\chapter{Introduction}
\label{ch:introduction}

\section{Motivation: two failures that look unrelated}

Modern machine learning exhibits two persistent, well-documented failure classes that
are almost always studied in isolation. The first is \emph{compositional generalisation}
failure: a model trained on known components and combinations fails to recombine them
under novel distribution, performing dramatically worse on combinations that are
structurally trivial yet unseen during training \citep{lake2018generalization,
keysers2020measuring,hupkes2023taxonomy}. The second is \emph{alignment} failure: a
system optimises competently toward objectives that diverge from what its designers
intended, whether through misspecification \citep{amodei2016concrete}, learned
optimisation pressures \citep{hubinger2019risks}, or catastrophic risk profiles
\citep{hendrycks2023overview}.

Both failure classes share a striking signature. In compositional generalisation
failure, the model's internal representations have collapsed onto the surface
statistics of the training distribution; the model has, in effect, lost the structural
content that would let it compose known parts \citep{lake2018generalization,
geirhos2020shortcut}. In alignment failure, the system's internal objective has
drifted from the specification it was trained to satisfy; the system pursues a proxy
that is coherent \emph{internally} but detached from the intended semantics
\citep{amodei2016concrete,russell2019human}. In both cases, what has failed is the
relationship between the internal structure of the system and the structure of the
world it is supposed to track.

\section{The thesis: a cumulative claim hierarchy}

This monograph develops its claim in three tiers, deliberately structured so that each
tier can stand or fall on its own evidence.

\paragraph{Primary thesis.}
\begin{quote}
\emph{Schema coherence is a measurable property of learned representation structure
whose deterioration can produce a stable low-coherence regime (the $\sigma$-trap)
associated with systematic generalisation failure.}
\end{quote}

\paragraph{Secondary thesis (extension).}
\begin{quote}
\emph{The same mechanism provides a testable account of some safety-relevant forms of
objective divergence} --- in particular where internal optimisation develops around
representations that have become decoupled from the intended semantic structure.
\end{quote}

\paragraph{Long-term implication (bounded).}
\begin{quote}
\emph{If these relationships persist in more capable systems, schema-coherent training
may constitute a candidate component of alignment strategy.}
\end{quote}

We call the internal degree of consistency of an agent's representation structure its
\emph{schema coherence} $\sigma$ (Chapter~4 gives the operational definition). The
primary thesis asserts that under gradient-based learning, $\sigma$ behaves like an
order parameter of a dynamical system: above a critical value $\sigmacrit$ learning is
compositional and robust; below it, the system is trapped in a stable equilibrium of
low coherence from which ordinary training does not escape. The $\Sigma$-Model, a
minimal mechanistic model developed in Chapter~7, is designed to exhibit exactly this
bifurcation \citep{amsyar2026sigma}.

It is important to be precise about what the thesis does \emph{not} claim. A stronger
formulation --- \emph{compositional generalisation failure and AI alignment failure are
the same phenomenon} --- has been considered and is explicitly rejected as the working
claim. Two complex, multiply-realised phenomena need not be identical as observables
for a common mechanism to connect them; claiming identity would invite the obvious
objection that a supervised benchmark failure and a sociotechnical alignment failure
are plainly different events. Instead, the thesis claims a \emph{unified-framework
hypothesis}: $\sigma$ is a common structural variable across both domains, and a
\emph{mechanistic hypothesis}: the $\sigma$-trap produces compositional generalisation
failure and, under the conditions specified in Chapter~8, safety-relevant optimisation
failures. Every link in the chain from compositional generalisation to alignment is
presented with an explicit evidential status --- literature-supported, observed,
modelled, experimentally manipulated, inferred, or speculative --- so that the reader
can see at each step what the evidence does and does not carry.

The thesis is deliberately structured so that no tier of the claim is presupposed by
the evidence reviews. Chapters~2--3 map the literature and the evidence base using
schema-coherence vocabulary strictly as a heuristic lens, and report what is found
--- including the limited explicit intersection between the compositional
generalisation and alignment literatures --- without requiring the claims to hold.
A literature gap establishes \emph{novelty}, not \emph{truth}; the reviews establish the
baseline, and the claims are argued cumulatively from Chapter~4 onward.

\subsection{What would falsify this thesis}

The mechanism claim is falsified if $\sigma$ is measured, interventions demonstrably
raise $\sigma$, and the predicted outcome does not change. Concretely, any of the
following would falsify the primary thesis:

\begin{enumerate}
  \item a low-$\sigma$ regime observed with intact systematic generalisation;
  \item a mediation analysis (Ch.~5) in which raising $\sigma$ leaves compositional
    generalisation unchanged;
  \item a $\Sigma$-Model in which the transition to non-compositional learning is not
    predicted by an independently measured $\sigma$ (Ch.~7).
\end{enumerate}

The secondary thesis additionally concedes the converse direction explicitly: some
alignment failures may lie entirely outside the framework, and Chapter~9 states which
ones those are. A thesis without a falsifiability criterion reads as dogma; these
criteria are the price of the claim being scientific.

\section{Why a monograph}

The claim spans two research communities that rarely read each other: the
compositional generalisation community, which works on benchmarks and mechanistic
accounts of systematicity \citep{keysers2020measuring,hupkes2023taxonomy,
lake2023compositional}, and the AI safety community, which works on specification,
robustness, and existential risk \citep{bostrom2014superintelligence,
hendrycks2023overview,ji2023survey}. A unified claim needs a unified presentation:
definitions that serve both communities, a shared notation, and a continuous argument
rather than a collection of standalone papers. The monograph form makes that possible.
Underlying publications --- the scoping review, the systematic review, the empirical
$\Sigma$-Model study, and the papers that will follow --- are listed in the front
matter and archived under \texttt{thesis/publications/}; chapters adapt them into
monograph prose rather than reproducing venue formatting.

\section{Contributions}

The monograph makes the following contributions.

\begin{enumerate}
  \item \textbf{Diagnosis.} A scoping review of the AGI safety landscape (Ch.~2) and a
    systematic review of the $\sigma$-trap evidence base (Ch.~3) establish that the
    safety literature is fragmented and that no existing framework treats internal
    schema coherence as the central failure mode.
  \item \textbf{Framework.} The $\Sigma$-Align framework (Ch.~4) formalises schema
    coherence as a measurable, intervenable quantity and gives definitions for the
    $\sigma$-trap and its bifurcation.
  \item \textbf{Evidence.} A pilot study of $\sigma$-coupling interventions (Ch.~5),
    a meta-analysis quantifying the $\sigma$-trap (Ch.~6), the $\Sigma$-Model showing
    compositional generalisation failure as a bifurcation (Ch.~7), and a study linking
    the trap to mesa-optimization (Ch.~8).
  \item \textbf{Implications.} A bounded argument that schema-coherent training is a
    candidate component of alignment strategy, with limitations stated before
    implications, and with direct relevance to corrigibility, CEV, and indirect
    normativity framed as long-term programme questions rather than established
    results (Ch.~9).
\end{enumerate}

\section{Reader's guide}

Chapter~2 maps the landscape of AGI safety research and identifies the gaps a
schema-coherence framework could address. Chapter~3 reviews the evidence for the
$\sigma$-trap as a measurable construct, study by study. Chapter~4 gives the
operational definition of $\sigma$ and its construct-validity tests. Chapters~5--8
present the empirical programme: the pilot (testing the causal chain
intervention $\rightarrow \sigma \rightarrow$ generalisation), the meta-analysis,
the $\Sigma$-Model (the flagship mechanism test), and the safety-relevant
optimisation study. Chapter~9 asks what the evidence can and cannot justify
concluding about alignment, and Chapter~10 concludes.

The chapters are written to be read in order, but each chapter opens with a summary of
its position in the argument and closes with a ``Relation to Other Chapters'' section,
so that readers from either community can enter at the chapter most relevant to them.
